\documentclass[UTF8]{ctexart}

\usepackage[a4paper,margin=1in]{geometry}
\usepackage{amsmath,amssymb,mathtools,bm}
\usepackage{booktabs,longtable,array}
\usepackage{hyperref}
\usepackage{enumitem}

\hypersetup{
    colorlinks=true,
    linkcolor=black,
    urlcolor=blue,
    citecolor=black
}

\setlist[itemize]{leftmargin=2em}
\setlist[enumerate]{leftmargin=2em}

\title{《Multirate Factor Analysis Models for Fault Detection in Multirate Processes》方法报告}
\author{}
\date{2026-03-27}

\begin{document}

\maketitle
\tableofcontents
\newpage

\section{文档目标}

本文围绕 Zhou 等人在 2019 年提出的 MR-FA（Multirate Factor Analysis，多采样率因子分析）方法，给出一份面向工程实现的中文技术报告。报告重点回答四个问题：
\begin{itemize}
    \item 论文到底提出了什么模型；
    \item 这个模型的输入、输出和中间量分别是什么，形状如何；
    \item 训练和推理时的数据流是怎样流动的；
    \item 相关数学原理应该如何理解，尤其是共享潜因子、EM 训练和故障检测统计量。
\end{itemize}

需要先说明一点：这篇论文本身不是深度学习论文，原文使用的是矩阵和向量记号，而不是神经网络里常见的“张量层”。因此，本文会在\textbf{保留论文原意}的基础上，把它\textbf{改写成工程上更容易实现的张量视角}。

\section{方法总览}

\subsection{一句话概括}

MR-FA 的核心思想是：\textbf{把不同采样率的变量都看成由同一组潜在公共因子驱动}，只是不同采样率在不同时间点上是否观测到、每组变量的维数、噪声强度都不一样。模型通过共享潜因子把“同采样率内部相关性”和“不同采样率之间的相关性”统一起来，而不是先上采样/下采样再做普通 FA。

\subsection{论文试图解决的问题}

工业过程里常见三类变量：
\begin{itemize}
    \item 高频在线传感器变量；
    \item 中频过程变量；
    \item 低频实验室质量变量。
\end{itemize}

传统 MSPM 方法大多默认所有变量同频采样。常见补救方法是：
\begin{itemize}
    \item 下采样到最慢频率；
    \item 上采样并插补慢变量；
    \item 只用高频变量，忽略低频质量变量。
\end{itemize}

这些做法都会损失信息或者扭曲真实相关结构。MR-FA 的目标就是直接在多采样率观测上建模，而不强行把数据压成单一频率。

\section{从 FA 到 MR-FA}

\subsection{单速率 FA 回顾}

普通因子分析（FA）假设观测向量 $\bm{x} \in \mathbb{R}^{M}$ 来自一个低维潜变量 $\bm{t} \in \mathbb{R}^{D}$：
\begin{equation}
    \bm{x} = \bm{\Psi}\bm{t} + \bm{\varepsilon},
\end{equation}
其中：
\begin{itemize}
    \item $\bm{\Psi} \in \mathbb{R}^{M \times D}$ 是因子载荷矩阵；
    \item $\bm{t} \sim \mathcal{N}(\bm{0}, \bm{I})$；
    \item $\bm{\varepsilon} \sim \mathcal{N}(\bm{0}, \bm{\Omega})$；
    \item $\bm{\Omega}$ 是对角阵，而不是 PPCA 中的 $\sigma^2 \bm{I}$。
\end{itemize}

这意味着 FA 允许每个变量拥有不同噪声方差，因此比 PPCA 的各向同性噪声假设更灵活。

\subsection{MR-FA 的关键变化}

对多采样率情形，论文把单个 FA 扩展为 $N$ 个采样率组：
\begin{equation}
    \bm{x}_n(k) = \bm{\Psi}_n \bm{t}(k) + \bm{\xi}_n(k), \quad n = 1,2,\dots,N.
\end{equation}

这里最关键的不是“有 $N$ 个方程”，而是\textbf{所有采样率组共享同一个时刻的潜因子 $\bm{t}(k)$}。也就是说：
\begin{itemize}
    \item 第 $n$ 组变量有自己的载荷矩阵 $\bm{\Psi}_n$；
    \item 第 $n$ 组变量有自己的对角噪声协方差 $\bm{\Omega}_n$；
    \item 但它们在同一个全局时间索引 $k$ 上由同一潜状态 $\bm{t}(k)$ 驱动。
\end{itemize}

这使得跨采样率相关性不需要显式写成交叉协方差块，而是通过共享潜因子自然耦合起来。

\section{符号与张量形状}

\subsection{论文原始表示}

设整个系统包含 $N$ 个采样率组，第 $n$ 组数据矩阵为
\begin{equation}
    \bm{X}_n \in \mathbb{R}^{K_n \times M_n},
\end{equation}
其中：
\begin{itemize}
    \item $K_n$ 是第 $n$ 个采样率组实际采到的样本数；
    \item $M_n$ 是第 $n$ 个采样率组包含的变量个数。
\end{itemize}

论文再引入一个\textbf{全局时间轴}，长度记为 $K$，满足
\begin{equation}
    K \ge \max(K_1,\dots,K_N),
\end{equation}
并且每个全局时间点 $k$ 至少有一组变量被采到。

\subsection{核心符号表}

\begin{longtable}{>{\raggedright\arraybackslash}p{0.18\linewidth} >{\raggedright\arraybackslash}p{0.22\linewidth} >{\raggedright\arraybackslash}p{0.50\linewidth}}
\toprule
符号 & 形状 & 含义 \\
\midrule
\endhead
$N$ & 标量 & 采样率组个数 \\
$M_n$ & 标量 & 第 $n$ 组变量维度 \\
$K_n$ & 标量 & 第 $n$ 组实际采样次数 \\
$K$ & 标量 & 统一后的全局时间步数 \\
$D$ & 标量 & 公共潜因子维度 \\
$\bm{x}_n(k)$ & $\mathbb{R}^{M_n}$ & 第 $n$ 组在全局时刻 $k$ 的观测向量 \\
$\bm{X}_n$ & $\mathbb{R}^{K_n \times M_n}$ & 第 $n$ 组完整训练数据矩阵 \\
$\phi_{k,n}$ & $\{0,1\}$ & 采样指示器，表示时刻 $k$ 是否观测到第 $n$ 组 \\
$\bm{X}_o(k)$ & $\mathbb{R}^{m_k}$ & 时刻 $k$ 所有已观测组拼接后的向量，$m_k=\sum_n \phi_{k,n}M_n$ \\
$\bm{t}(k)$ & $\mathbb{R}^{D}$ & 时刻 $k$ 的公共潜因子 \\
$\hat{\bm{t}}(k)$ & $\mathbb{R}^{D}$ & 潜因子后验均值 \\
$\bm{\Psi}_n$ & $\mathbb{R}^{M_n \times D}$ & 第 $n$ 组载荷矩阵 \\
$\bm{\Omega}_n$ & $\mathbb{R}^{M_n \times M_n}$ & 第 $n$ 组噪声协方差，对角阵 \\
$\bm{\Sigma}(k)$ & $\mathbb{R}^{D \times D}$ & 论文符号中的后验精度矩阵 \\
$\mathbb{E}[\bm{t}(k)\bm{t}(k)^\top \mid \bm{X}_o(k)]$ & $\mathbb{R}^{D \times D}$ & 后验二阶矩 \\
$T^2(k)$ & 标量 & 公共因子子空间监控统计量 \\
$\mathrm{SPE}_n(k)$ & 标量 & 第 $n$ 组残差平方统计量 \\
\bottomrule
\end{longtable}

\subsection{工程实现视角下的张量化表示}

论文的原始输入并不是单个规整的三维张量，因为不同采样率组的维数 $M_n$ 往往不同。因此最自然的实现方式其实是：
\begin{equation}
    \{\bm{X}_1,\bm{X}_2,\dots,\bm{X}_N\},
\end{equation}
也就是一个“矩阵列表”。

如果为了向量化实现，也可以构造填充后的张量：
\begin{equation}
    \mathcal{X}_{\mathrm{pad}} \in \mathbb{R}^{K \times N \times M_{\max}}, \quad M_{\max}=\max_n M_n,
\end{equation}
并配套使用掩码：
\begin{equation}
    \mathcal{M}_{\mathrm{pad}} \in \{0,1\}^{K \times N \times M_{\max}},
\end{equation}
以及组级采样指示器：
\begin{equation}
    \bm{\Phi} \in \{0,1\}^{K \times N}.
\end{equation}

\paragraph{结论}
从论文角度看，\textbf{输入不是一个标准 dense tensor，而是多组不同形状矩阵加一个全局时间对齐机制}；从工程角度看，可以把它实现成“填充张量 + 掩码”或“按组存储 + 观测模式分桶”。

\section{模型输入、输出与中间结果}

\subsection{训练输入}

训练集由正常工况数据组成，输入包括：
\begin{itemize}
    \item 多采样率观测集合 $\{\bm{X}_n\}_{n=1}^{N}$；
    \item 全局时间轴上的采样指示器 $\phi_{k,n}$。
\end{itemize}

如果按统一时间轴表示，则单个时间步的输入是：
\begin{equation}
    \left\{\bm{x}_n(k), \phi_{k,n}\right\}_{n=1}^{N}.
\end{equation}
真正参与当前步推断的观测是拼接向量 $\bm{X}_o(k)$。

\subsection{训练输出}

训练结束后得到的不是类别标签，而是一个生成式监测模型：
\begin{itemize}
    \item 各组载荷矩阵 $\{\bm{\Psi}_n\}_{n=1}^{N}$；
    \item 各组噪声协方差 $\{\bm{\Omega}_n\}_{n=1}^{N}$；
    \item 训练数据上各时刻的后验因子均值 $\hat{\bm{t}}(k)$；
    \item 每个时刻的 $T^2$ 统计量和每个采样率组对应的 $\mathrm{SPE}_n$ 统计量；
    \item 基于训练分布估计出的控制限。
\end{itemize}

\subsection{测试时输出}

对新时刻 $k_{\text{new}}$ 的观测，模型输出：
\begin{itemize}
    \item 因子估计 $\hat{\bm{t}}(k_{\text{new}}) \in \mathbb{R}^{D}$；
    \item 公共子空间统计量 $T^2(k_{\text{new}})$；
    \item 每个采样率组的残差统计量 $\mathrm{SPE}_n(k_{\text{new}})$。
\end{itemize}

如果批量处理 $K_{\text{test}}$ 个测试时刻，那么：
\begin{equation}
    \hat{\bm{T}} \in \mathbb{R}^{K_{\text{test}} \times D}, \quad
    \bm{t}^2 \in \mathbb{R}^{K_{\text{test}}}, \quad
    \mathrm{SPE} \in \mathbb{R}^{K_{\text{test}} \times N}.
\end{equation}

\section{MR-FA 的数学结构}

\subsection{生成模型}

MR-FA 在每个时间步上是一个线性高斯潜变量模型。对第 $n$ 个采样率组：
\begin{equation}
    \bm{x}_n(k) = \bm{\Psi}_n \bm{t}(k) + \bm{\xi}_n(k),
\end{equation}
其中
\begin{align}
    \bm{t}(k) &\sim \mathcal{N}(\bm{0}, \bm{I}), \\
    \bm{\xi}_n(k) &\sim \mathcal{N}(\bm{0}, \bm{\Omega}_n), \\
    \bm{\Omega}_n &= \operatorname{diag}(\omega_{n,1},\dots,\omega_{n,M_n}).
\end{align}

这意味着：
\begin{itemize}
    \item 潜因子之间先验不相关；
    \item 同组变量的噪声在论文里被假设为相互独立；
    \item 不同组之间的统计耦合不是通过噪声项给出，而是通过共享潜因子给出。
\end{itemize}

\subsection{采样指示器的作用}

由于某些组在时刻 $k$ 没有采到，所以后验推断不能直接把所有组硬拼起来。论文用
\begin{equation}
    \phi_{k,n} =
    \begin{cases}
        1, & \text{时刻 } k \text{ 观测到第 } n \text{ 组}, \\
        0, & \text{否则}
    \end{cases}
\end{equation}
来决定当前时刻有哪些组参与推断。

因此，MR-FA 是一个\textbf{结构随时间变化的模型}：时刻 $k$ 真正生效的并不是全部参数块，而是所有满足 $\phi_{k,n}=1$ 的那部分参数块。这也是论文把 MR-FA 视为一种特殊时变模型的原因。

\subsection{后验推断}

固定当前参数 $\{\bm{\Psi}_n,\bm{\Omega}_n\}$ 时，时刻 $k$ 的潜因子后验均值为：
\begin{equation}
    \hat{\bm{t}}(k)
    =
    \bm{\Sigma}(k)^{-1}
    \sum_{n=1}^{N}
    \phi_{k,n}\bm{\Psi}_n^\top \bm{\Omega}_n^{-1}\bm{x}_n(k).
\end{equation}

其中
\begin{equation}
    \bm{\Sigma}(k)
    =
    \sum_{n=1}^{N}
    \phi_{k,n}\bm{\Psi}_n^\top \bm{\Omega}_n^{-1}\bm{\Psi}_n
    + \bm{I}.
\end{equation}

后验二阶矩为：
\begin{equation}
    \mathbb{E}\!\left[\bm{t}(k)\bm{t}(k)^\top \mid \bm{X}_o(k)\right]
    =
    \bm{\Sigma}(k)^{-1}
    + \hat{\bm{t}}(k)\hat{\bm{t}}(k)^\top.
\end{equation}

\paragraph{一个容易混淆但很重要的点}
论文把 $\bm{\Sigma}(k)$ 记成了一个类似“协方差”的符号，但从公式结构看，它实际上是\textbf{后验精度矩阵}，因为真正的后验协方差是 $\bm{\Sigma}(k)^{-1}$。这一点在阅读公式和做代码实现时必须分清。

\section{张量形状分析}

\subsection{单时间步上的形状传播}

对某个时刻 $k$，第 $n$ 组的关键量形状如下：
\begin{align}
    \bm{x}_n(k) &: M_n \times 1, \\
    \bm{\Psi}_n &: M_n \times D, \\
    \bm{\Omega}_n^{-1} &: M_n \times M_n, \\
    \bm{\Psi}_n^\top \bm{\Omega}_n^{-1} \bm{x}_n(k) &: D \times 1.
\end{align}

于是多组贡献求和后得到：
\begin{equation}
    \sum_{n=1}^{N}\phi_{k,n}\bm{\Psi}_n^\top \bm{\Omega}_n^{-1} \bm{x}_n(k)
    \in \mathbb{R}^{D}.
\end{equation}

再乘上 $\bm{\Sigma}(k)^{-1}\in\mathbb{R}^{D\times D}$，得到
\begin{equation}
    \hat{\bm{t}}(k)\in\mathbb{R}^{D}.
\end{equation}

\subsection{M 步中的形状传播}

固定第 $n$ 组，更新载荷矩阵时：
\begin{equation}
    \hat{\bm{\Psi}}_n
    =
    \left(
        \sum_{k=1}^{K}
        \phi_{k,n}\bm{x}_n(k)\hat{\bm{t}}(k)^\top
    \right)
    \left(
        \sum_{k=1}^{K}
        \phi_{k,n}
        \mathbb{E}\!\left[\bm{t}(k)\bm{t}(k)^\top \mid \bm{X}_o(k)\right]
    \right)^{-1}.
\end{equation}

其中：
\begin{itemize}
    \item $\bm{x}_n(k)\hat{\bm{t}}(k)^\top$ 的形状是 $M_n \times D$；
    \item 分母是 $D \times D$；
    \item 所以最终 $\hat{\bm{\Psi}}_n$ 形状仍是 $M_n \times D$。
\end{itemize}

噪声协方差更新为：
\begin{equation}
    \hat{\bm{\Omega}}_n
    =
    \operatorname{diag}
    \left(
        \bm{S}_n
        -
        \frac{1}{K_n}
        \sum_{k=1}^{K}
        \phi_{k,n}\hat{\bm{\Psi}}_n \hat{\bm{t}}(k)\bm{x}_n(k)^\top
    \right),
\end{equation}
其中
\begin{equation}
    \bm{S}_n = \frac{1}{K_n}\sum_{k=1}^{K_n}\bm{x}_n(k)\bm{x}_n(k)^\top \in \mathbb{R}^{M_n \times M_n}.
\end{equation}

\subsection{批量实现时最常见的张量}

若按工程代码批量实现，一般会出现下列张量：
\begin{itemize}
    \item 观测模式矩阵 $\bm{\Phi}\in\{0,1\}^{K\times N}$；
    \item 后验均值矩阵 $\hat{\bm{T}}\in\mathbb{R}^{K\times D}$；
    \item 后验精度张量 $\mathcal{P}\in\mathbb{R}^{K\times D\times D}$；
    \item 后验协方差张量 $\mathcal{C}\in\mathbb{R}^{K\times D\times D}$；
    \item 二阶矩张量 $\mathcal{S}\in\mathbb{R}^{K\times D\times D}$；
    \item 监测统计量 $\bm{t}^2\in\mathbb{R}^{K}$；
    \item 残差统计量矩阵 $\mathrm{SPE}\in\mathbb{R}^{K\times N}$。
\end{itemize}

\paragraph{对应到仓库实现}
仓库中的 \texttt{baseline/mr\_fa.py} 基本就是这个张量化思路：
\begin{itemize}
    \item \texttt{phi}: $(K,N)$；
    \item \texttt{means}: $(K,D)$；
    \item \texttt{precisions}: $(K,D,D)$；
    \item \texttt{covariances}: $(K,D,D)$；
    \item \texttt{second\_moments}: $(K,D,D)$；
    \item \texttt{spe}: $(K,N)$。
\end{itemize}

\section{训练数据流：从输入到参数更新}

\subsection{总流程}

MR-FA 的训练可以写成下面的数据流：
\begin{equation}
    \{\bm{X}_n\}_{n=1}^{N}
    \longrightarrow
    \{\bm{X}_o(k), \phi_{k,:}\}_{k=1}^{K}
    \xrightarrow{\text{E-step}}
    \{\hat{\bm{t}}(k), \mathbb{E}[\bm{t}(k)\bm{t}(k)^\top]\}_{k=1}^{K}
    \xrightarrow{\text{M-step}}
    \{\bm{\Psi}_n,\bm{\Omega}_n\}_{n=1}^{N}.
\end{equation}

训练只使用正常工况数据，因此本质上是一个无监督建模过程。

\subsection{E 步：估计隐因子后验}

在 E 步，参数固定，逐时刻推断公共因子。可以把它理解为：
\begin{enumerate}
    \item 先看当前时刻有哪些组被观测到；
    \item 把这些组各自对潜因子的“证据”映射到 $D$ 维潜空间；
    \item 把所有证据相加；
    \item 再乘以后验协方差，得到当前时刻的最佳因子估计。
\end{enumerate}

由于模型是线性高斯的，上述结果是标准高斯共轭推断的闭式解。

\subsection{M 步：更新载荷和噪声}

在 M 步，潜因子后验矩固定，用最大似然更新参数：
\begin{itemize}
    \item $\bm{\Psi}_n$ 由“观测与潜因子”的交叉矩估计；
    \item $\bm{\Omega}_n$ 由该组的残差协方差对角部分估计。
\end{itemize}

注意这里虽然更新是按组进行的，但\textbf{每个组的更新并不是只用本组数据}。原因在于：
\begin{itemize}
    \item $\hat{\bm{t}}(k)$ 是由当前时刻\textbf{所有已观测组}共同决定的；
    \item 所以一旦第 $n$ 组在时刻 $k$ 出现，它的参数更新就间接受到了其他组观测的影响。
\end{itemize}

这正是论文强调的“MR-FA 不等价于多个独立 FA”的根本原因。

\subsection{为什么这能建模跨采样率相关性}

如果为每个采样率单独训练一个 FA：
\begin{itemize}
    \item 每组只看到自己；
    \item 每组有自己的潜因子；
    \item 组间没有共享隐状态，因此无法利用跨组共现信息。
\end{itemize}

MR-FA 则相反：
\begin{itemize}
    \item 各组有各自的观测头 $\bm{\Psi}_n$；
    \item 但共享同一潜空间 $\bm{t}(k)$；
    \item 只要多个采样率在某些时间点共现，它们就会通过同一个 $\bm{t}(k)$ 彼此约束。
\end{itemize}

\section{推理与故障检测数据流}

\subsection{测试时的因子推断}

对新时刻 $k_{\text{new}}$ 的观测，仍然使用与训练 E 步相同的后验推断公式：
\begin{equation}
    \hat{\bm{t}}(k_{\text{new}})
    =
    \bm{\Sigma}(k_{\text{new}})^{-1}
    \sum_{n=1}^{N}
    \phi_{k_{\text{new}},n}
    \bm{\Psi}_n^\top \bm{\Omega}_n^{-1}\bm{x}_n(k_{\text{new}}).
\end{equation}

\subsection{公共子空间统计量 $T^2$}

论文定义：
\begin{equation}
    T^2(k_{\text{new}})
    =
    \hat{\bm{t}}(k_{\text{new}})^\top
    \bm{\Sigma}(k_{\text{new}})^{-1}
    \hat{\bm{t}}(k_{\text{new}}).
\end{equation}

直观上，$T^2$ 衡量的是“当前样本在公共潜空间中偏离正常模式的程度”。如果一个故障会显著改变多个采样率共享的潜在机理，那么它往往会首先在 $T^2$ 上表现出来。

\subsection{残差统计量 $\mathrm{SPE}$}

对每个采样率组，先做重构：
\begin{equation}
    \hat{\bm{x}}_n(k) = \bm{\Psi}_n \hat{\bm{t}}(k),
\end{equation}
再计算残差：
\begin{equation}
    \bm{e}_n(k) = \bm{x}_n(k) - \hat{\bm{x}}_n(k),
\end{equation}
最后得到
\begin{equation}
    \mathrm{SPE}_n(k) = \bm{e}_n(k)^\top \bm{e}_n(k).
\end{equation}

$\mathrm{SPE}_n$ 反映的是“第 $n$ 组变量中无法被共享潜因子解释的那部分能量”。因此：
\begin{itemize}
    \item $T^2$ 更偏向监控共享子空间；
    \item $\mathrm{SPE}_n$ 更偏向监控各组残差子空间。
\end{itemize}

\subsection{控制限}

论文采用：
\begin{itemize}
    \item $T^2$ 的控制限由 $F$ 分布估计；
    \item $\mathrm{SPE}_n$ 的控制限由卡方分布估计。
\end{itemize}

因此完整检测流程可以写成：
\begin{equation}
    \text{新观测}
    \rightarrow
    \hat{\bm{t}}
    \rightarrow
    \left(T^2,\mathrm{SPE}_1,\dots,\mathrm{SPE}_N\right)
    \rightarrow
    \text{与控制限比较}
    \rightarrow
    \text{报警/不报警}.
\end{equation}

\section{几何解释}

\subsection{单个 FA 的几何意义}

单个 FA 可以理解为：在 $M$ 维观测空间中寻找一个 $D$ 维超平面，让样本投影到这个平面后，尽可能保留原始变量之间的相关结构。观测向量可以分解为：
\begin{equation}
    \bm{x} = \hat{\bm{x}} + \bm{\xi},
\end{equation}
其中 $\hat{\bm{x}}$ 是投影部分，$\bm{\xi}$ 是残差。

\subsection{MR-FA 的几何意义}

MR-FA 不是分别给每个采样率找一个超平面，而是让\textbf{所有采样率组共享同一个潜在超平面}。因此它优化的不是某一组内部的局部误差最小，而是所有采样率组总体残差协方差的全局最小。

论文的几何解释可以概括为：
\begin{itemize}
    \item 先只用一种采样率时，会得到一个超平面；
    \item 再加入另一种采样率后，最优超平面会旋转；
    \item 这个旋转的结果，是为了让所有采样率组在同一个共享潜空间上达到更好的整体拟合。
\end{itemize}

这个观点很重要，因为它解释了为什么 MR-FA 与“多个独立 FA 模型”在本质上不同。

\section{与其他方法的差异}

\subsection{与下采样/上采样方法的差异}

MR-FA 不先把所有变量压成同频数据，因此避免了两个问题：
\begin{itemize}
    \item 下采样导致高频信息损失；
    \item 上采样导致人工插值破坏真实相关性。
\end{itemize}

\subsection{与多个独立 FA 的差异}

独立 FA 的问题是每组潜因子互不共享，因此：
\begin{itemize}
    \item 不能显式学习跨采样率相关性；
    \item 低频质量变量无法反向约束高频过程变量的潜空间；
    \item 最终更像多个局部模型，而不是一个统一过程模型。
\end{itemize}

\subsection{与 MRPCA 的差异}

论文把 MR-FA 与之前的 MRPCA 进行了比较。最关键的差异是：
\begin{itemize}
    \item MRPCA 仍然继承了 PPCA 风格的各向同性噪声假设；
    \item MR-FA 允许每个变量有自己的噪声方差，因此对工业数据更现实。
\end{itemize}

对工业过程数据而言，不同变量的噪声幅值差异通常很大，所以这一点确实重要。

\section{方法的优点、假设与局限}

\subsection{优点}

\begin{itemize}
    \item 不需要先统一采样率；
    \item 共享潜因子，能同时利用组内和组间相关性；
    \item 比 MRPCA 更灵活，因为噪声方差不是各向同性；
    \item 输出既有潜空间统计量 $T^2$，也有分组残差统计量 $\mathrm{SPE}_n$，监测解释性较好。
\end{itemize}

\subsection{核心假设}

\begin{itemize}
    \item 线性生成关系：$\bm{x}_n=\bm{\Psi}_n\bm{t}+\bm{\xi}_n$；
    \item 潜因子服从标准高斯先验；
    \item 每组内部噪声协方差为对角阵；
    \item 不同采样率的耦合主要通过共享潜因子，而不是通过动态时序模型显式建模。
\end{itemize}

\subsection{局限}

\begin{itemize}
    \item 这是静态线性模型，不直接建模时间动态；
    \item 如果不同采样率互为互素并且共现很少，跨组耦合会变弱；
    \item 对角噪声协方差仍然是一种简化，无法表示同组变量残差之间的相关噪声；
    \item EM 只能保证收敛到局部最优。
\end{itemize}

\section{面向实现的数据流总结}

\subsection{训练期}

用更接近代码的方式，训练流程可以概括为：
\begin{enumerate}
    \item 读入各采样率组数据 $\{\bm{X}_n\}$；
    \item 建立全局时间轴，并生成组级观测掩码 $\bm{\Phi}$；
    \item 初始化 $\{\bm{\Psi}_n,\bm{\Omega}_n\}$；
    \item 对每个观测模式重复执行 E 步，得到所有时刻的 $\hat{\bm{t}}(k)$ 和二阶矩；
    \item 对每个采样率组执行 M 步，更新 $\bm{\Psi}_n,\bm{\Omega}_n$；
    \item 重复 4--5 直到对数似然收敛；
    \item 在训练集上计算 $T^2$ 和 $\mathrm{SPE}_n$ 的控制限。
\end{enumerate}

\subsection{推理期}

\begin{enumerate}
    \item 输入新时刻的多采样率观测；
    \item 根据当前观测模式计算 $\hat{\bm{t}}$；
    \item 对每个已观测组做重构并求残差；
    \item 输出 $T^2$ 和各组 $\mathrm{SPE}_n$；
    \item 与控制限比较，给出报警结果。
\end{enumerate}

\subsection{如果一定要写成“输入张量 $\rightarrow$ 输出张量”}

可以写成：
\begin{align}
    \mathcal{X}_{\mathrm{pad}} &\in \mathbb{R}^{K \times N \times M_{\max}}, \\
    \bm{\Phi} &\in \{0,1\}^{K \times N}, \\
    \hat{\bm{T}} &\in \mathbb{R}^{K \times D}, \\
    \mathcal{C} &\in \mathbb{R}^{K \times D \times D}, \\
    \mathrm{SPE} &\in \mathbb{R}^{K \times N}.
\end{align}

但要强调：\textbf{这只是工程上的统一封装方式，不是论文最原生的数学表达}。论文原生表达更接近“多组矩阵 + 条件观测模式”。

\section{论文实验结论简述}

论文用两个案例验证方法：
\begin{itemize}
    \item Tennessee--Eastman 模拟过程；
    \item 造纸废水处理中的 R2S 厌氧反应器真实数据。
\end{itemize}

整体结论是：
\begin{itemize}
    \item MR-FA 普遍优于下采样 FA、只用高频变量的 FA，以及早期的 MRPCA；
    \item 尤其在需要依赖低频质量变量才能识别的故障上，MR-FA 的优势更明显；
    \item 论文将性能优势主要归因于两点：完整利用多采样率信息，以及更合理的异方差噪声建模。
\end{itemize}

\section{结论}

MR-FA 可以看成是\textbf{多采样率场景下的共享潜因子生成模型}。它的本质不是补齐缺失数据，也不是给每种采样率单独建模，而是在统一时间轴上用一个共享潜空间解释所有可见观测。对工程实现而言，最重要的理解有三点：
\begin{itemize}
    \item 输入本质上是“多组不同形状矩阵 + 采样指示器”，不是天然规则的单张量；
    \item 模型的核心输出是各时刻公共因子后验、各组载荷矩阵与噪声方差，以及由此得到的 $T^2/\mathrm{SPE}$ 统计量；
    \item 整个数据流由“观测模式 $\rightarrow$ 后验因子 $\rightarrow$ 参数更新/残差统计”三段组成，其中跨采样率信息融合发生在共享潜因子的后验推断阶段。
\end{itemize}

\end{document}
